% GENERATED from content/data/enrichment.json by mxbuild.enrichment; do not edit.
\mxunit{enrichment-telescoping-notes}{Sommes télescopiques}
\mxsubtitle{Première--Terminale · suites \textperiodcentered{} Leçon et méthode \textperiodcentered{} 10}

\begin{prerequisites}
Fractions, sommes finies et indices.
\end{prerequisites}

\begin{keypoint}[Idées et méthode]
Une somme télescopique s’écrit comme une somme de différences consécutives. Les termes intérieurs s’annulent : $\sum (a_{k}-a_{k+1})=a_{1}-a_{n+1}$ pour k allant de 1 à n.

Écrire d’abord trois termes et les deux derniers permet de voir les annulations et de contrôler les bornes. Une simplification terme à terme ne dispense pas de préciser les indices.

Une décomposition utile est $\frac{1}{k(k+1)}=\frac{1}{k}-\frac{1}{k+1}$, valable pour $k\geq 1$. On la vérifie en mettant les fractions au même dénominateur.
\end{keypoint}

\begin{example}
$\sum  \frac{1}{k(k+1)}$, de $k=1$ à n, vaut $1-\frac{1}{n+1}=\frac{n}{n+1}$. La somme est inférieure à 1 et tend vers 1 lorsque n augmente.
\end{example}

\begin{activity}[À explorer]
Calculer les cinq premières sommes partielles, puis étudier l’écart à 1. Distinguer la somme finie de sa limite.
\end{activity}

